% !PRINTED 513
\weberpage{513}

Ersetzt man in (2) $v$ durch eine andere Klasseninvariante derselben Diskriminante $4D$ und desselben Geschlechtes, so geht auch $u$ in eine andere Klasseninvariante der Diskriminante $D$ über, bleibt aber auch in demselben Geschlecht.

Läßt man also in (2) $v$ die Klasseninvarianten \gesperrt{eines Geschlechtes der Diskriminante} $4D$ durchlaufen, so durchläuft $u$ die Invarianten des entsprechenden Geschlechtes der Diskriminante $D$ [indem es jeden dieser Werte dreimal, bei $D\equiv 1\pmod 8$ nur einmal annimmt], und die symmetrischen Funktionen dieser $u$ sind in dem gleichen Rationalitätsbereich enthalten wie die Größen $v$. Ist daher die Klassenfunktion $H_{-4D}(v)$ nach den Geschlechtern in Faktoren zerlegbar, so gilt das gleiche von $H_{-D}(u)^{1)}$.

\textbf{§ 138. Zerfällung der Klassengleichung nach den Geschlechtern.}

Um den Satz 1., § 136, anzuwenden, setze man
\weq{1}{D=B^2-4AC=-4m,}
und zerlege $m$ in zwei Faktoren
\weq{2}{m=m'm'',}
\gesperrt{wobei vorausgesetzt ist, daß $m''$ ungerade und ohne quadratischen Teiler sei.}

Nun machen wir in dem Satz 1., § 136, die Annahme:
\weq{3}{\begin{aligned}
x&=4,\quad y=2(4m'-m''),\\
n&=(4m'+m'')^2.
\end{aligned}}
\weq{4}{\begin{aligned}
\alpha&=2B+4m'-m'',\\
\alpha m''&=4AC-(B-m'')^2,\\
\beta&=4A.
\end{aligned}}
\weq{5}{\begin{aligned}
\alpha+\beta\omega&=2i\sqrt{m'm''}+4m'-m'',\\
&=(2\sqrt{m'}+i\sqrt{m''})^2,
\end{aligned}}
% SOURCE-ORIGINAL-SIC: Equation (5) prints 2i sqrt(m'm'') in its first line, while the square in its second line expands with 4i sqrt(m'm''); the printed coefficient 2 is preserved.

\footnotetext[1]{Die Durchführung der entsprechenden Betrachtungen für ein ungerades $D$ würde zwar auch möglich sein, würde aber zahlreiche Unterscheidungen und Weitläufigkeiten nötig machen. Hier ist ein Punkt, wo die \gesperrt{Gauss}sche Bezeichnung und Unterscheidung von Formen erster und zweiter Art, deren ich mich noch in der ersten Auflage bedient habe, eine gewisse Vereinfachung des Ausdrucks mit sich bringen würde. Es hängt das mit der Ausnahmestellung zusammen, die die Zahl 2 in der ganzen Theorie der elliptischen Funktionen einnimmt, die in der \gesperrt{Weierstrass}schen Theorie etwas zurücktritt, aber doch nicht ganz verschwindet. Auf der anderen Seite ist diese Auszeichnung der Zahl 2 auch wieder die Quelle von großen Vereinfachungen, namentlich in den numerischen Resultaten.}
%CONT>>
